\section{Technische Anforderungen}
Die technischen Anforderungen beschreiben die notwendigen Eigenschaften und Randbedingungen für die Umsetzung des Systems.
Sie werden in die Bereiche Hardware, Software, Schnittstellen und Performance unterteilt.

\subsection{Hardware}
Die Hardwarearchitektur basiert auf einem Mikrocontroller mit integrierter \gls{wlan}-Schnittstelle.
Dieser übernimmt sowohl die Datenverarbeitung als auch die Kommunikation mit externen Diensten.
Die Anzeige erfolgt über mehrere \gls{led}-Streifen mit \gls{ws2812b} \gls{rgb}-Leuchtdiode.
Eine \gls{led}-Streifen wird für die Hauptanzeige verwenet. Das ist die Uhrzeit in Stunden und Minuten, die Außentemperautur und das Wetter.
Der zweite \gls{led}-Streifen wird für die Sekundenanzeige verwendet.
Der dritte \gls{led}-Streifen für die Anzeige der Tageszeit, der Raumfeuchtigkeit, der Mondphase und der Innenraumtemperatur.
Alle Hardware Anforderungen sind in Tabelle~\ref{tab:Hardware Anforderung} aufgelistet.

\begin{table}[H]
\centering
\begin{tabular}{p{3cm} p{9cm}}
\textbf{ID} & \textbf{Beschreibung} \\
\hline
H1 & Verwendung eines Mikrocontrollers mit WLAN-Funktionalität \\
H2 & Ansteuerung von einzeln adressierbaren LEDs \\
H3 & 3 Kaskdierte LED Strukturen mit je einem Datenpin am Mikrocontroller \\
H4 & Einbindung eines Temperatur- und Feuchtigkeitssensors \\
H5 & Sicherstellung einer stabilen Spannungsversorgung für alle Komponenten \\
\end{tabular}
\caption{Technische Anforderungen – Hardware}
\label{tab:Hardware Anforderung}
\end{table}

\subsection{Software}
Die Softwarearchitektur ist modular aufgebaut.
Zentrale Komponenten sind die Zeitverarbeitung, die \gls{led}-Ansteuerung sowie die Kommunikation mit externen Diensten.
Die Verarbeitung erfolgt regelbasiert und ermöglicht eine flexible Erweiterung.
Alle Software Anforderungen sind in Tabelle~\ref{tab:Software Anforderung} aufgelistet.

\begin{table}[H]
\centering
\begin{tabular}{p{3cm} p{9cm}}
\textbf{ID} & \textbf{Beschreibung} \\
\hline
S1 & Implementierung einer Zeitverwaltung mit NTP-Synchronisation \\
S2 & Regelbasierte Umwandlung numerischer Zeitwerte in Wortdarstellung \\
S3 & Verarbeitung und Klassifikation von Wetterdaten \\
S4 & Implementierung einer Sensorabfrage für Innenraumwerte \\
S5 & Steuerung der LED-Anzeige in Echtzeit \\
S6 & Bereitstellung eines HTTP-Servers zur Benutzerinteraktion \\
S7 & Fehlerbehandlung bei Kommunikations- und Sensorfehlern \\
\end{tabular}
\caption{Technische Anforderungen – Software}
\label{tab:Software Anforderung}
\end{table}

\subsection{Schnittstellen}
Die Schnittstellen ermöglichen die Integration externer Datenquellen sowie die Interaktion mit dem System.
Die Kommunikation erfolgt überwiegend über standardisierte Protokolle.
Alle Schnittstellen Anforderungen sind in Tabelle~\ref{tab:Schnittstellen Anforderung} aufgelistet.

\begin{table}[H]
\centering
\begin{tabular}{p{3cm} p{9cm}}
\textbf{ID} & \textbf{Beschreibung} \\
\hline
I1 & WLAN-Schnittstelle zur Internetanbindung \\
I2 & HTTP-Schnittstelle zur Benutzerinteraktion \\
I3 & API-Anbindung zur Abfrage von Wetterdaten \\
I4 & Sensorschnittstelle zur Erfassung von Temperatur und Luftfeuchtigkeit \\
\end{tabular}
\caption{Technische Anforderungen – Schnittstellen}
\label{tab:Schnittstellen Anforderung}
\end{table}

\subsection{Performance}
Die Performanceanforderungen stellen sicher, dass das System eine flüssige und aktuelle Anzeige gewährleistet.
Gleichzeitig werden externe Daten in definierten Intervallen aktualisiert, um eine ausreichende Aktualität zu erreichen.
Alle Schnittstellen Anforderungen sind in Tabelle~\ref{tab:Performance Anforderung} aufgelistet.

\begin{table}[H]
\centering
\begin{tabular}{p{3cm} p{9cm}}
\textbf{ID} & \textbf{Beschreibung} \\
\hline
P1 & Aktualisierung der Anzeige in Intervallen kleiner gleich 1 s \\
P2 & Wetterdatenaktualisierung mindestens alle 15 min \\
P3 & Sensoraktualisierung im Bereich weniger Sekunden \\
P4 & Minimierung von Verzögerungen bei der Benutzerinteraktion \\
\end{tabular}
\caption{Technische Anforderungen – Performance}
\label{tab:Performance Anforderung}
\end{table}

% %%%%%%%%%%%Old Version
% Die technischen Anforderungen legen Randbedingungen für die Umsetzung auf Hardware- und Softwareebene fest.

% \subsection*{Hardware}

% Die Hardware soll eine homogene und ausreichend helle Ausleuchtung der Wortmatrix ermöglichen. Die eingesetzten LEDs sollen einzeln oder gruppenweise ansteuerbar sein. Die maximale Leistungsaufnahme soll 30~W nicht überschreiten.

% Ein Mikrocontroller mit ausreichender Rechenleistung und integrierter Netzwerkschnittstelle soll die zentrale Steuerung übernehmen. Der verfügbare Flash- und RAM-Speicher soll die Implementierung der Anzeige- und Kommunikationsfunktionen ermöglichen.

% Ein Echtzeituhr-Modul oder eine netzwerkbasierte Zeitsynchronisation soll eine dauerhafte Zeitgenauigkeit sicherstellen. Ein Helligkeitssensor soll die Umgebungsbeleuchtung erfassen.

% \subsection*{Software}

% Die Softwarearchitektur soll modular aufgebaut sein. Die Module für Zeitberechnung, Anzeige, Kommunikation und Sensorik sollen logisch getrennt implementiert werden.

% Die Firmware soll eine stabile WLAN-Verbindung unterstützen. Die Zeitsynchronisation soll über ein standardisiertes Protokoll erfolgen. Die Berechnung der Mondphase soll auf einem geeigneten astronomischen Modell basieren.

% Die Software soll Updatefähigkeit vorsehen. Eine Konfigurationsoberfläche soll die Anpassung zentraler Parameter ermöglichen.

% \subsection*{Schnittstellen}

% Das System soll mindestens eine drahtlose Netzwerkschnittstelle nach IEEE 802.11 bereitstellen. Zusätzlich soll eine serielle Schnittstelle für Entwicklungs- und Diagnosezwecke vorgesehen werden.

% Sensoren und Anzeigeelemente sollen über standardisierte digitale Schnittstellen angebunden werden. Die Versorgung soll über ein externes Netzteil mit geeigneter Spannungsstabilisierung erfolgen.

% \subsection*{Performance}

% Die Aktualisierung der Anzeige soll innerhalb von 500~ms erfolgen. Die Reaktionszeit auf Benutzereingaben soll unter 1~s liegen.

% Die Systemstabilität soll einen Dauerbetrieb von mindestens 24~h ohne Neustart ermöglichen. Die mittlere Leistungsaufnahme im Normalbetrieb soll 15~W nicht überschreiten.

% \subsection{Optische Anforderungen} \todo{text}
% Die optischen Anforderungen definieren die gestalterischen Ziele der Wortuhr.

% Das Design soll eine klare und reduzierte Formensprache aufweisen. Die Wortmatrix soll gleichmäßig ausgeleuchtet sein. Einzelne Lichtpunkte außerhalb aktiver Wörter sollen nicht sichtbar sein.

% Die Schriftart soll gut lesbar sein. Der Kontrast zwischen Hintergrund und aktiven Wörtern soll ausreichend hoch sein. Die Anzeige soll aus einem Betrachtungsabstand von 3~m eindeutig erkennbar sein.

% Das Gehäuse soll sich in typische Wohnräume integrieren lassen. Die sichtbaren Materialien sollen hochwertig und langlebig wirken. Kabel und Befestigungselemente sollen von außen nicht sichtbar sein.
